\noindent
Eine Funktion heißt \textbf{gerade}, wenn sie für alle $x$ aus ihrem Definitionsbereich
die folgende Eigenschaft erfüllt:
\begin{equation}
f(-x) = f(x)
\end{equation}

\noindent
Der Graph einer geraden Funktion ist achsensymmetrisch zur y-Achse. Das bedeutet,
dass der Funktionswert für $x$ und $-x$ identisch ist.

\noindent
Eine Funktion heißt \textbf{ungerade}, wenn für alle $x$ gilt:
\begin{equation}
f(-x) = -f(x)
\end{equation}

\noindent
Ungerade Funktionen sind punktsymmetrisch zum Ursprung. Der Funktionswert an der
Stelle $-x$ ist also das negative des Funktionswertes an der Stelle $x$